\subsection{Finite rank cells on the continuous spectral shrinkage path}
\label{sec:v25-bernstein-rank-cells}
The V23 interval method bounded source-trial score signs by
enclosing the individual rational terms. In low-dimensional
EEG, a complementary approach acts on their \emph{common
numerator}, exploiting cancellations otherwise lost in
termwise bounds.

For fixed source-trained eigensystem
$S=U\,\mathrm{diag}(s_1,\ldots,s_p)U^\top$, class contrast
$d_s$, fixed held-\emph{source} positive--negative feature
difference $v_{pn}$, let
\[
q_i(\lambda)=(1-\lambda)s_i+\lambda\nu,\qquad
a_{i,pn}=(u_i^\top v_{pn})(u_i^\top d_s).
\]
On a feasible interval $I$ where $q_i(\lambda)>0$ for all
$i,\lambda$, the sign of the score difference is the sign
of the single polynomial
\begin{equation}
 N_{pn}(\lambda)
 =\sum_{i=1}^{p}a_{i,pn}\prod_{j\ne i}q_j(\lambda),
 \qquad\deg N_{pn}\le p-1 .
 \label{eq:v25-spectral-numerator}
\end{equation}
Indeed $z_{pn}=N_{pn}/\prod_iq_i$ and the denominator is
strictly positive. Unless $N_{pn}$ vanishes identically
(and the pair remains tied throughout), a positive--negative
pair can change its strict ordering only at the at most
$p-1$ distinct real roots of $N_{pn}$. Hence the finite
source-validation macro-AUC is piecewise constant on
the complement of the finite union of all such rank-cell
boundaries. Values \emph{at} a root are evaluated separately
because ties count for one-half AUC credit; it would be
incorrect to assume the supremum is attained only at
open-cell midpoints. The root-count argument is an
elementary consequence of rational functions, not a
claim of priority for new polynomial algebra.

For each subinterval $[\ell,r]$, map $\lambda$ to
$t\in[0,1]$ and express the degree-$p-1$ numerator
as a Bernstein polynomial
\[
 N_{pn}(\ell+(r-\ell)t)
 =\sum_{k=0}^{p-1} b_{k,pn}B_k^{p-1}(t).
\]
As Bernstein basis functions are nonnegative and sum to
one, $\min_k b_{k,pn}\le N_{pn}(\lambda)\le
\max_k b_{k,pn}$. Strict positivity (or negativity)
of \emph{all} coefficients proves that the pair's rank
is unchanged on the entire subinterval. A source pair
is declared rank-stable over $I$ only if \emph{every}
subinterval proves the same strict sign. A failure to
certify is retained as uncertainty, not counted as a
real crossing. The inference is exact in real arithmetic;
our floating-point implementation applies a conservative
cancellation margin but does not use formally verified
directed rounding, and therefore we do not claim a
machine-checked mathematical certificate.

The new implementation uses a positive product-of-denominators
scaling and is explicitly restricted to $p\le12$ to
control numerical conditioning. With three EEG channels,
a tangent-space covariance model has only six unique
coordinates, yielding a numerator of degree at most
five per positive--negative sample pair. AUC credit of
all certified source pairs remains unchanged over the
entire candidate path. Adding one-sided worst-case
credit changes for the uncertified pairs provides a
deterministic interval-wide finite-\emph{source}-AUC
envelope analogous to Eq.~\eqref{eq:v23-one-sided-auc-upper}.

\paragraph{Research scope and falsification.}
The practical hypothesis is that whole-numerator
Bernstein bounds reduce unnecessary uncertainty relative
to the V23 termwise/Taylor interval method on
\emph{identical frozen} held-source model folds and
lambda intervals. We test this on a previously exposed
three-channel Weibo2014 acquisition as a retrospective
mechanistic comparison, not as a fresh target-validation
sample. Either method may be tighter; gains in certified
\emph{source} rank coverage are not gains in unknown-person
AUC, and no claim of cross-subject error prevention
is implied. Boundary roots, exact ties, denser intervals,
numerical cancellation, and subject-specific harms remain
necessary negative controls.
